\section{Thin cuboids}\label{sec:thin}
For side lengths at least two, write $F(a,b,c)$ for the least inventory
size guaranteeing a frame, and $G(a,b,c)$ for the corresponding complete-box
threshold. These extend the cube notation above. A frame contains
$4(a+b+c)-16$ positions; the target colouring may depend on the inventory.

\begin{theorem}[Rods and strips]\label{thm:thin}
Put $K_\ell=\max\{29,\lfloor(24\ell-1)/5\rfloor\}$. Then
\[
 F(2,2,\ell)=G(2,2,\ell)=K_\ell \qquad(\ell\ge2),
\]
and
\[
 F(2,3,\ell)=K_\ell,\qquad G(2,3,\ell)=\max\{29,6\ell\}
 \qquad(\ell\ge3).
\]
\end{theorem}
Despite their four additional positions, the strip frames have the same
threshold as the rods. Every rod can require spare sources, with the ratio
of threshold to volume tending to $6/5$. In contrast, every strip of length
at least five can be assembled from exactly its volume: any thirty sources,
for example, form a $2\times3\times5$ box.

We prove one upper bound for both families. Deleting the four unit-edge
positions from a $2\times3\times\ell$ frame leaves the requirements of a
$2\times2\times\ell$ rod, with the same corner and long-edge placements.
The case $\ell=2$ is the corner theorem. All subsequent upper-bound arguments
therefore concern strips with $\ell\ge3$.

\paragraph{Lower bounds.}
The twenty-eight-cube obstruction of Theorem~\ref{thm:sharp} prevents the
corners of either shape. For the other bound, put
\[
 N_0=\left\lfloor\frac{24\ell-6}{5}\right\rfloor
     =\left\lfloor\frac{24\ell-1}{5}\right\rfloor-1
\]
and distribute $N_0$ sources as evenly as possible among the six
checkerboards. Each bichromatic type is missing from at least
$\lfloor N_0/6\rfloor$ sources. The inequality
$5N_0\le24\ell-6$ gives
\[
 \left\lfloor N_0/6\right\rfloor\ge N_0-4\ell+1,
\]
so every type has at most $4\ell-1$ suppliers.
All source corners are trichromatic. Hence any possible three-colour target
must have equal-coloured opposite faces, of type $D$: faces of the same
colour cannot be adjacent, since that would give a bicoloured corner.
For every orientation of $D$, its eight corners and $4(\ell-2)$ long-edge
positions all require the same bichromatic pair. These $4\ell$ positions
cannot be supplied. This proves both frame lower bounds, regardless of the
four extra strip positions.

\paragraph{Choosing directions for an opposite-face target.}
Reduce $N=K_\ell$ sources by corner containment. As in
Section~\ref{sec:frames}, this also preserves every edge requirement.
Put $q=\ell-2$ and $s=N-4\ell$. Directly from the definition,
\begin{equation}\label{eq:thin-bounds}
 N\ge29,\qquad s\ge5,\qquad N\le6s+5.
\end{equation}
In particular $N\ge4\ell+4$, the number of strip-frame positions.
Index the six checkerboards by their missing pairs, and let $x_e$ be their
multiplicities, viewed as edge weights on $K_4$. The incident weight
$z_v=\sum_{e\ni v}x_e$ counts checkerboards supporting both orientations
of the triple omitting colour $v$. Call $v$ good if that triple admits a
$D$ corner solution. In particular, $z_v\ge8$ implies that $v$ is good.

\begin{lemma}\label{lem:thin-pairs}
If some vertex is good, there are a good vertex $v$ and two distinct edges
$e,f$ of its complementary triangle such that
\[
 x_e\le s,\qquad x_f\le N-12.
\]
\end{lemma}
\begin{proof}
First suppose every edge weight is at most $N-12$. If no good vertex has a
complementary edge of weight at most $s$, every complementary edge of any
good vertex has weight at least $s+1$. Starting with one good vertex, each
other vertex then has incident weight at least $2(s+1)\ge12$ and is good.
Every edge is complementary to some good vertex, forcing total weight at
least $6s+6>N$, a contradiction. Choose a light complementary edge $e$;
any other edge of that triangle can serve as $f$.

Otherwise some edge has weight at least $N-11$. Either endpoint $v$ is good,
since its incident weight is at least $N-11\ge18$. The other five edges
have total weight at most eleven. Thus the triangle complementary to $v$
has an edge of weight at most three and all its edges have weight at most
eleven. These satisfy the required bounds by~\eqref{eq:thin-bounds}.
\end{proof}

Write $e=ab$ and $f=ad$. Orient the target as $(a,a,b,b,d,d)$,
with axis lengths $(2,3,\ell)$. Its noncorner demands
are $4q$ copies of $ab$ and four copies of $ad$. Fix a corner solution.
Every one of its eight sources possesses both pairs, so the remaining
supplies are at least $N-x_{ab}-8\ge4q$ and
$N-x_{ad}-8\ge4$. Every reduced source lacks at most one bichromatic pair,
so the two supplier sets together contain every remaining source.
There are at least $4q+4$ such sources. Hall's theorem completes the frame.
Thus any $D$ corner solution suffices.

\paragraph{When no opposite-face target has a corner solution.}
Let $A\ge B$ count the two ordinary orientation classes, and $C$ the
checkerboards. Since $z_v\le7$ and $\sum_v z_v=2C$, we have $C\le14$.
Consequently $A\ge8$. If $B+z_v\ge4$, four of those sources and four
majority sources supply the two orientations of a $D$ target, a
contradiction. Hence
\begin{equation}\label{eq:thin-minority}
 B+z_v\le3\quad\hbox{for every }v,\qquad
 4B+2C\le12,\qquad B+C\le6,\qquad A\ge23.
\end{equation}
In particular each checkerboard multiplicity is at most three.

We will reserve corners and any required monochromatic edges first.
If $r$ sources have been reserved, each bichromatic type has at least
$N-r-3$ remaining suppliers. Any two distinct bichromatic types together
see all $N-r$ remaining sources. Therefore Hall's inequalities hold as soon
as every individual demand is at most $N-r-3$ and the total demand is at
most $N-r$. This applies to any number of distinct bichromatic types:
a position set meeting at least two types sees all remaining sources.

If a majority bucket contains eight cubes, Lemma~\ref{lem:bucket} gives
corners of a $Q(a;b,d)$ target. Orient it as $(b,d,a,a,a,a)$.
Its monochromatic edges run along the length-two direction and have no
noncorner positions. The first row of Table~\ref{tab:thin-demands} now
completes the frame.

Otherwise every majority bucket has at most seven cubes. Write $t_a$ for
its $T_a$ count and $n_T,n_P$ for the majority totals. Then
\[
 n_T+2n_P\le28,\qquad n_T+n_P=A\ge23,\qquad n_T\ge18.
\]
Every $t_a\le7$, and at most one colour has $t_a<2$: two such colours
would give $n_T\le1+1+7+7=16$. There is a minority source, since otherwise
$N=A\ge29$ contradicts the bucket bound. Each minority source supports a
triple containing any specified colour $a$. An ordinary minority source
supplies the opposite orientation; a checkerboard supplies both.

For distinct $a,b,d$, let $H$ have corner multiset
\[
 2aab+2aad+2abb+abd+adb.
\]
Four majority $T_a$ sources and two majority $T_b$ sources fill the six
bicoloured positions; the minority source and one further majority source
fill the two trichromatic positions. Two orientations of $H$ give the
last two rows of the following table. Face tuples use the order
$(+x,-x,+y,-y,+z,-z)$, with axis lengths $(2,3,\ell)$.

\begin{table}[ht]
\centering
\begin{tabular}{llllr}
\toprule
Target & Face tuple & Monochromatic demand & Bichromatic demands & $r$\\
\midrule
$Q$ & $(b,d,a,a,a,a)$ & none & $ab:2q+2,\ ad:2q+2$ & 8\\
$H$ & $(a,a,b,d,a,b)$ & $aa:2$ & $ab:2q+2,\ ad:2q$ & 10\\
$H$ & $(b,d,a,a,a,b)$ & $bb:1$ & $ab:2q+1,\ ad:2q+1,\ bd:1$ & 9\\
\bottomrule
\end{tabular}
\caption{Noncorner strip demands; $r$ sources supply the corners and
monochromatic edges before the remaining Hall matching.}
\label{tab:thin-demands}
\end{table}

If some $t_a\ge6$, choose a minority-supported triple containing $a$ and
choose another colour $b$ in it with $t_b\ge2$. Such a colour exists by
the preceding bound. Reserve six $T_a$, two $T_b$, the minority cube and
one unused majority cube. These ten distinct sources fill the corners
and the two $aa$ positions in the second row.

If no $t_a\ge6$, then every $t_a\le5$. The sum is at least eighteen,
so every $t_a\ge3$ and some $t_a\ge5$. Choose such an $a$, and either
other colour $b$ of a minority-supported triple. Four $T_a$, three $T_b$
and the two trichromatic sources fill the corners and the $bb$ position
in the third row. There are ample unused majority sources for the
trichromatic assignment in both constructions, since $A\ge23$.

For each row the total bichromatic demand is $4\ell+4-r\le N-r$.
The largest demand is at most $N-r-3$: for $\ell=3$ this follows from
$N\ge29$, and for $\ell\ge4$ from $N\ge4\ell+5$.
The preceding Hall observation completes every construction. This proves
the strip upper bound, and hence both frame formulas.

\paragraph{Completion and shape dependence.}
A rod consists entirely of its frame. In a strip, every position outside
the frame exposes just one exterior face, whose colour can be supplied by
any unused source. The complete-box threshold is therefore the maximum of
volume and frame threshold. Since $\lfloor(24\ell-1)/5\rfloor<6\ell$,
this proves the claimed formula for $G(2,3,\ell)$.
The threshold depends on shape, not just frame cardinality:
$2\times2\times7$ and $2\times3\times6$ both have twenty-eight frame
positions, but their thresholds are respectively thirty-three and
twenty-nine.
